\documentclass{estilo}

\usepackage[spanish]{babel}
\usepackage{graphicx}
\usepackage{float}
\usepackage{amsmath}        % para los vectores columnas
\usepackage{amsfonts}       % para las negrita de pizarra
\usepackage{amssymb}        % para simbolos matematicos
\usepackage{multirow}       % para las tablas
%\usepackage{hyperref} 
\usepackage{dsfont}
\usepackage{listings}
\usepackage{xcolor}
\usepackage{enumitem,multicol,setspace} %por las dudovich
\usepackage{ragged2e}       % para que no tire error con \justifying
\usepackage{tikz} % para generar las imagenes bien piola sin berretin
\usepackage{caption}
\usetikzlibrary{arrows.meta} % tikz⁺

% estilos globales acá en el main para que afecten a todo
\tikzset{
  vertex/.style={circle, draw=blue!70!black, fill=blue!10, very thick, minimum size=9mm, font=\sffamily\bfseries},
  dag edge/.style={-{Stealth[length=2.4mm,width=1.7mm]}, semithick, draw=gray!75!black, shorten >=2pt, shorten <=2pt}
}





\definecolor{codegreen}{rgb}{0,0.6,0}
\definecolor{codegray}{rgb}{0.5,0.5,0.5}
\definecolor{codepurple}{rgb}{0.58,0,0.82}
\definecolor{backcolour}{rgb}{0.95,0.95,0.92}
\definecolor{codeblue}{rgb}{0.16,0.36,0.67}
% config. para que el pseudocodigo tenga colorcitos.
\lstset{
    backgroundcolor=\color{backcolor},   % color de fondo, ese lo ponemos?
    commentstyle=\color{codegreen},       % color de comentarios
    keywordstyle=\color{codepurple}\bfseries, % color de palabras clave (if, return, def)
    numberstyle=\tiny\color{codegray}, % color de números de línea
    stringstyle=\color{codeblue}, % color de strings
    basicstyle=\ttfamily\small, % tipografia monoespaciada
    breakatwhitespace=false,         
    breaklines=true,                 
    captionpos=b,                    
    keepspaces=true,                 
    numbers=right,
    numbersep=5pt,                  
    showspaces=false,                
    showstringspaces=false,
    showtabs=false,                  
    tabsize=4,
    language=Python % detecta automáticamente def, return, if, else, etc.
}
\lstdefinestyle{mystyle}{
    backgroundcolor=\color{backcolour},   
    commentstyle=\color{codegreen},
    keywordstyle=\color{magenta},
    numberstyle=\tiny\color{codegray},
    stringstyle=\color{codepurple},
    basicstyle=\ttfamily\footnotesize,
    breakatwhitespace=false,         
    breaklines=true,                 
    captionpos=b,                    
    keepspaces=true,                 
    numbers=left,                    
    numbersep=5pt,                  
    showspaces=false,                
    showstringspaces=false,
    showtabs=false,                  
    tabsize=2
}
\lstset{style=mystyle}

\usepackage{enumitem,multicol,setspace}
\onehalfspacing % Interlineado 1.5
\setlength{\parskip}{1.5ex plus 0.5ex minus 0.2ex} % Espacio entre párrafos
\setlength{\parindent}{0pt} % Elimina la sangría obligatoria para que no quede raro 


\usepackage{hyperref} 
\hypersetup{
    hidelinks,
    pdfstartview={XYZ null null null},
    pdfview={XYZ null null null}
}

\usepackage[style=authoryear-comp]{biblatex}
\addbibresource{referencias.bib}


\newcounter{subenum}[enumi] % para las multicolumnas

\renewcommand{\thesubenum}{\arabic{subenum}}
\usepackage[nomessages]{fp}
\FPeval\thecolwidth{round(1/4:4)}% Specify number of columns -> column width
\newcommand{\newitem}[1]{%
  \refstepcounter{subenum}%
  \parbox{\dimexpr\thecolwidth\linewidth-.5\columnsep}{%
    \makebox[\labelwidth][r]{(\thesubenum)\hspace*{\labelsep}}%
    #1}\hfill%
}


\usepackage{scalerel,stackengine} % para el sombrero
\stackMath
\newcommand\rhat[1]{%
\savestack{\tmpbox}{\stretchto{%
  \scaleto{%
    \scalerel*[\widthof{\ensuremath{#1}}]{\kern-.6pt\bigwedge\kern-.6pt}%
    {\rule[-\textheight/2]{1ex}{\textheight}}%WIDTH-LIMITED BIG WEDGE
  }{\textheight}% 
}{0.5ex}}%
\stackon[1pt]{#1}{\tmpbox}%
}
\parskip 1ex

\usepackage{mathtools}      % floor y ceil
\DeclarePairedDelimiter\ceil{\lceil}{\rceil}
\DeclarePairedDelimiter\floor{\lfloor}{\rfloor} 




\begin{document}
\nocite{*} % para la bibliografía.
\maketitle

\justifying

\newpage
\begin{center}
    \vspace*{0.1cm} 
    \Huge\bfseries Índice
    \vspace{0.1cm}
\end{center}

\renewcommand{\contentsname}{} % elimina el indice chico que mete latex por defecto

% El comando automático ahora va a renderizar links cliqueables a tus archivos
\tableofcontents

\clearpage 
%\section{Consigna 1: Selección de Intervalos}
%\subsection{Enfoque Greedy}
% aparece en el indice pero sin numero: \addcontentsline{toc}{section}{Conclusion}


\newpage
\vspace*{1.5cm}
\section{Resumen}

Se presenta un resumen del Trabajo Práctico 3, detallando los temas a tratar para ofrecer una lectura fluida, clara y precisa que permita comprender el desarrollo en su totalidad.

En primer lugar, se va a hablar sobre el problema propuesto, cómo fue el planteo de la resolución, cómo se logró llegar a que ésta misma pueda ser óptima para cualquier caso posible, las consecuencias temporales que genera eso y cómo empeorar la optimalidad nos puede traer una gran ventaja.

Después, se trata con profundidad el análisis de la complejidad de los problemas y cómo es posible categorizarlosrizarlos en conjuntos siempre y cuando cumplan condiciones polinómicas y/o exponenciales.

Se detalla la teoría necesaria sobre los algoritmos de \textbf{Backtracking}, algoritmos de \textbf{Programación Lineal} y algoritmos de \textbf{Aproximación}. Para ver estos aspectos aplicados, se incluye un breve ejemplo de la idea en general y su analogía a conceptos en el mundo real (toma de decisiones en la vida). Por otro lado, se van a detallar el diseño de los algoritmos implementados en \textit{pseudocódigo}, explicando anteriormente la lógica, cómo fue buscada y cómo fue pensada esa resolución. Adjuntándose al final la implementación realizada.

Se realiza un \textit{análisis} de cada algoritmo planteado, mediante el estudio de la complejidad y una demostración matemática para validar la optimalidad o cota del algoritmo, de forma clara y con ejemplos propuestos. Se incorporan múltiples experimentos empíricos de forma tal que se pueda visualizar gráficamente las medidas de tiempo sobre cierta cantidad de casos y poder comparar todos los algoritmos implementados.

Finalmente, se expone una conclusión sobre lo aprendido, la dificultad del problema, y la particularidad de estos algoritmos y su pensamiento.

\newpage
\section{El Problema}
Los hermanos siguieron creciendo. Mateo también aprendió sobre programación dinámica, y cada uno aplicaba la lógica sabiendo que el otro también lo hacía. El juego de las monedas se tornó aburrido en cuanto notaron que siempre ganaba quien empezara, o según la suerte. Los años pasaron, llegó la adolescencia y empezaron a tener gustos diferentes. En general, jugaban a juegos individuales. En particular, Sophia estaba muy enganchada con el siguiente juego de tipo \textit{puzzle}: El juego consiste de una grilla de $n \times m$ celdas. Cada celda puede estar vacía, estar ocupada por una piedra verde, o estar ocupada por una piedra celeste. El objetivo original del juego es remover algunas de las piedras de tal forma que las piedras que quedan cumplan las siguientes condiciones:

\begin{enumerate}
    \item Cada fila tiene al menos una piedra.
    \item Ninguna columna tiene piedras de ambos colores.
\end{enumerate}

A Sophia le interesa tanto esta versión, como también aquella que \textbf{relaja} la primera restricción, obteniendo una solución que maximice la cantidad de fila que cumplen la restricción (manteniendo siempre la segunda restricción). Es decir, minimizar la cantidad de filas que quedan vacías.


Para transformar esta situación en un problema algorítmico concreto, el desafío es el siguiente:
Hacer un análisis del problema, y cumplir varios puntos (después explicados) del problema:
Dado un tablero de $n \times m$ celdas, donde cada celda está vacía u ocupada con una piedra verde o celeste (solo puede estar ocupada por 1 como mucho), ¿es posible remover piedras de tal forma que cada fila tenga al menos una piedra y cada columna no tenga piedras de ambos colores?

Los puntos a resolver son:
\begin{itemize}
    \item Demostrar que el Problema del juego de las Piedras se encuentra en NP, y es un problema \textit{NP-Completo}.
    \item Escribir para obtener la solución óptima al problema en su versión de optimización: un algoritmo de backtracking y un modelo de programación lineal
    \item Analizar el famoso algoritmo de aproximación de ``El Gringo Hinz'' (\ref{sec:hinz})
    \item Implementar un algoritmo de aproximación para el problema.
\end{itemize}
\subsection{Introducción}

Tranquilamente se pueden observar varios puntos a realizar. Lo ideal, como primera instancia, hay que tener en cuenta que existen tableros \textbf{sin} solución, los cuales no se pueden cumplir las dos condiciones iniciales del problema. A eso sumémosle que si se demuestra que este problema es NP-Completo, estamos frente a un problema \textit{muy difícil} de resolver.
Nos podemos preguntar:
\begin{itemize}
    \item ¿Es realmente posible encontrar una solución óptima? Si tanto se habla de que es muy difícil, ¿\textit{cuánto tardaría}?
    \item ¿Cuáles son las diferencias entre un algoritmo de Backtracking y Programación Lineal? ¿Distan mucho de los algoritmos de Aproximación? ¿Todos me aseguran optimalidad?
    \item ¿Cuánto me beneficia el Problema del Juego de las Piedras en su \textit{versión de Relajación}?
\end{itemize}

Una vez propuesto eso, sigamos pensando para ahora empezar con el planteo.

\subsection{Planteamiento y Modelado del problema}

Para cualquier tablero que se nos presente, nuestro objetivo va a ser ganar siempre que podamos.
En la Figura \ref{fig:con_solucion} se observa una instancia válida para el problema de optimización

\begin{figure}[H]
    \centering
    % tablero 3x3 con solución en versión de decisión
\begin{tikzpicture}[
    scale=0.85,
    transform shape,
    celda/.style={draw=gray!60, fill=gray!10, line width=0.6pt},
    verde/.style={circle, fill=green!60!black, draw=green!40!black, minimum size=0.68cm, inner sep=0pt},
    celeste/.style={circle, fill=cyan!75!blue!80, draw=cyan!50!black, minimum size=0.68cm, inner sep=0pt},
    tachado/.style={circle, draw=red!80!black, line width=1.3pt, densely dashed, minimum size=0.72cm, inner sep=0pt}
]

% macro local para dibujar la cruz roja sobre una piedra
\newcommand{\cruz}[2]{
    \node[tachado] at (#1, #2) {};
    \draw[red!80!black, line width=1.4pt] (#1-0.22, #2-0.22) -- (#1+0.22, #2+0.22);
    \draw[red!80!black, line width=1.4pt] (#1-0.22, #2+0.22) -- (#1+0.22, #2-0.22);
}

%  TABLERO INICIAL
\begin{scope}[shift={(0,0)}]
    \draw[celda] (0,0) rectangle (3,3);
    \draw[celda] (0,0) grid (3,3);
    
    % Fila 1 (arriba: y=2.5)
    \node[verde] at (0.5, 2.5) {};
    \node[celeste] at (1.5, 2.5) {};
    % Fila 2 (medio: y=1.5)
    \node[verde] at (0.5, 1.5) {};
    \node[celeste] at (2.5, 1.5) {};
    % Fila 3 (abajo: y=0.5)
    \node[celeste] at (1.5, 0.5) {};
    \node[verde] at (2.5, 0.5) {};
    
    \node[font=\footnotesize\sffamily] at (1.5, -0.45) {Tablero original};
\end{scope}

%  PIEDRAS A REMOVER 
\begin{scope}[shift={(4.2,0)}]
    \draw[celda] (0,0) rectangle (3,3);
    \draw[celda] (0,0) grid (3,3);
    
    % Fila 1
    \node[verde] at (0.5, 2.5) {};
    \node[celeste] at (1.5, 2.5) {}; \cruz{1.5}{2.5}
    % Fila 2
    \node[verde] at (0.5, 1.5) {}; \cruz{0.5}{1.5}
    \node[celeste] at (2.5, 1.5) {};
    % Fila 3
    \node[celeste] at (1.5, 0.5) {};
    \node[verde] at (2.5, 0.5) {}; \cruz{2.5}{0.5}
    
    \node[font=\footnotesize\sffamily] at (1.5, -0.45) {Piedras removidas};
\end{scope}

% RESULTADO VÁLIDO 
\begin{scope}[shift={(8.4,0)}]
    \draw[celda] (0,0) rectangle (3,3);
    \draw[celda] (0,0) grid (3,3);
    
    % Fila 1 (Col 1: Verde)
    \node[verde] at (0.5, 2.5) {};
    % Fila 2 (Col 3: Celeste)
    \node[celeste] at (2.5, 1.5) {};
    % Fila 3 (Col 2: Celeste)
    \node[celeste] at (1.5, 0.5) {};
    
    \node[font=\footnotesize\sffamily] at (1.5, -0.45) {Solución válida};
\end{scope}

\end{tikzpicture}

    \renewcommand{\thefigure}{1}
    \caption{Ejemplo de tablero $3 \times 3$ con solución en la versión de decision.}
    \label{fig:con_solucion}
\end{figure}

Por otro lado, la Figura \ref{fig:sin_solucion} muestra una configuración donde resulta imposible cubrir todas las filas sin generar conflictos de color en las columnas.

\begin{figure}[H]
    \centering
    %La fila 2 obliga a que la columna 1 sea celeste, y la fila 3 obliga a que la columna 2 sea verde. Con esa configuración, en la fila 1 deben eliminarse ambas piedras y queda vacía.

\begin{tikzpicture}[
    scale=0.85,
    transform shape,
    celda/.style={draw=gray!60, fill=gray!10, line width=0.6pt},
    verde/.style={circle, fill=green!60!black, draw=green!40!black, minimum size=0.7cm, inner sep=0pt},
    celeste/.style={circle, fill=cyan!75!blue!80, draw=cyan!50!black, minimum size=0.7cm, inner sep=0pt}
]

\draw[celda] (0,0) rectangle (3,3);
\draw[celda] (0,0) grid (3,3);

% Fila 1 (arriba: y=2.5) -> Verde en col 1, Celeste en col 2
\node[verde] at (0.5, 2.5) {};
\node[celeste] at (1.5, 2.5) {};

% Fila 2 (medio: y=1.5) -> Solo Celeste en col 1 (fuerza col 1 a ser celeste)
\node[celeste] at (0.5, 1.5) {};

% Fila 3 (abajo: y=0.5) -> Solo Verde en col 2 (fuerza col 2 a ser verde)
\node[verde] at (1.5, 0.5) {};

\end{tikzpicture}

    \renewcommand{\thefigure}{1.1}
    \caption{Ejemplo de tablero $3 \times 3$ sin solución para el problema.}
    \label{fig:sin_solucion}
\end{figure}

Finalmente, al relajar la primera restricción (Figura~\ref{fig:relajacion}), una instancia que no admite solución exacta logra maximizar la cantidad de filas cubiertas dejando únicamente una fila vacía.

\begin{figure}[H]
    \centering
    % sin solucion en decision -> optimo en relajacion
\begin{tikzpicture}[
    scale=0.85,
    transform shape,
    celda/.style={draw=gray!60, fill=gray!10, line width=0.6pt},
    verde/.style={circle, fill=green!60!black, draw=green!40!black, minimum size=0.65cm, inner sep=0pt},
    celeste/.style={circle, fill=cyan!75!blue!80, draw=cyan!50!black, minimum size=0.65cm, inner sep=0pt},
    tachado/.style={circle, draw=red!80!black, line width=1.3pt, densely dashed, minimum size=0.7cm, inner sep=0pt}
]

\newcommand{\cruz}[2]{
    \node[tachado] at (#1, #2) {};
    \draw[red!80!black, line width=1.4pt] (#1-0.21, #2-0.21) -- (#1+0.21, #2+0.21);
    \draw[red!80!black, line width=1.4pt] (#1-0.21, #2+0.21) -- (#1+0.21, #2-0.21);
}

% INSTANCIA INFACTIBLE EN DECISIÓN 
\begin{scope}[shift={(0,0)}]
    \draw[celda] (0,0) rectangle (3,4);
    \draw[celda] (0,0) grid (3,4);
    
    % Fila 1 (y=3.5): Verde en col 1, Celeste en col 2
    \node[verde] at (0.5, 3.5) {};
    \node[celeste] at (1.5, 3.5) {};
    % Fila 2 (y=2.5): Solo Celeste en col 1
    \node[celeste] at (0.5, 2.5) {};
    % Fila 3 (y=1.5): Solo Verde en col 2
    \node[verde] at (1.5, 1.5) {};
    % Fila 4 (y=0.5): Verde en col 3
    \node[verde] at (2.5, 0.5) {};
    
    \node[font=\footnotesize\sffamily] at (1.5, -0.45) {Tablero original};
\end{scope}

%  Seleccion de piedras removidas ÓPTIMA
\begin{scope}[shift={(4.2,0)}]
    \draw[celda] (0,0) rectangle (3,4);
    \draw[celda] (0,0) grid (3,4);
    
    % Fila 1 (queda vacía para salvar filas 2 y 3)
    \node[verde] at (0.5, 3.5) {}; \cruz{0.5}{3.5}
    \node[celeste] at (1.5, 3.5) {}; \cruz{1.5}{3.5}
    % Fila 2
    \node[celeste] at (0.5, 2.5) {};
    % Fila 3
    \node[verde] at (1.5, 1.5) {};
    % Fila 4
    \node[verde] at (2.5, 0.5) {};
    
    \node[font=\footnotesize\sffamily] at (1.5, -0.45) {Removidas óptimas};
\end{scope}

% SOLUCIÓN RELAJADA (3/4 FILAS
\begin{scope}[shift={(8.4,0)}]
    \draw[celda] (0,0) rectangle (3,4);
    \draw[celda] (0,0) grid (3,4);
    
    % Fila 1 vacía (y=3.5) - resaltada suavemente en rojo
    \draw[fill=red!8, draw=red!40, dashed] (0,3) rectangle (3,4);
    
    % Fila 2
    \node[celeste] at (0.5, 2.5) {};
    % Fila 3
    \node[verde] at (1.5, 1.5) {};
    % Fila 4
    \node[verde] at (2.5, 0.5) {};
    
    \node[font=\footnotesize\sffamily] at (1.5, -0.45) {Óptimo ($3$ de $4$ filas)};
\end{scope}

\end{tikzpicture}


    \renewcommand{\thefigure}{1.2}
    \caption{Resolución óptima de una instancia del problema en la versión relajada.}
    \label{fig:relajacion}
\end{figure}
Habiendo encontrado un contraejemplo para que si o si tenemos que dejar al menos una fila vacía, surge la pregunta clave para modelar el problema:
¿cómo tomamos las decisiones sin perdernos en las $n \times m$ celdas?

Si analizamos la segunda regla (\textit{ninguna columna puede tener piedras de ambos colores}), notamos que: una vez que decidimos que en una columna $j$ no va a haber piedras celestes, nunca nos va a convenir eliminar ninguna piedra verde de esa misma columna, porque conservar todas las piedras verdes disponibles nunca rompe la restriccion vertical y al mismo tiempo esta maximizando nuestras chances de cubrir filas horizontales.
Por lo tanto, el problema de elegir qué piedras sacar es matemáticamente igual a agarrar exactamente \textbf{$m$ decisiones binarias} (una por cada columna $j \in \{0, \dots, m-1\}$): o la columna $j$ sigue solo con sus piedras verdes (\texttt{V}), o deja solo sus piedras celestes (\texttt{C}).

Por último, es muy importante diferenciar entre las dos variantes del problema:
\begin{itemize}
    \item \textbf{Versión de Decisión:} Dado un tablero de $n \times m$ celdas, ¿es posible remover piedras de forma tal que ninguna columna tenga piedras de ambos colores y todas las filas tengan \textbf{al menos} una piedra ($0$ filas vacías)? (¿es posible dejar a lo sumo $k$ filas vacías?).
    \item \textbf{Versión de Optimización (Relajada):} Dado un tablero de $n \times m$ celdas, encontrar una asignación válida de colores por columna que \textbf{minimice la cantidad de filas que quedan vacías} (o equivalentemente, que maximice la cantidad de filas cubiertas).
\end{itemize}

\subsection{Consecuencias de la optimalidad}

Si se logra pensar y crear estrategias para las cuales siempre que se pueda ganar lo logremos de forma eficaz, eso significa que:
\begin{itemize}
    \item Las indicaciones propuestas fueron óptimas.
    \item Si pasamos las instrucciones a un algoritmo, ese mismo es óptimo.
    \item Si la estrategia no devuelve siempre el óptimo, sino una solucion \textit{cercana} a la óptima, la solución se considera una \textit{aproximada}.
    \item Si logramos resolverlo polinomialmente, el algoritmo se considera rápido y muy eficaz si nos devuelve siempre la solución óptima. En cambio, si logramos resolverlo exponencialmente habría que considerar \textit{sacrificar} cierta optimalidad para que, con tableros de gran magnitud, no tardemos semanas en resolverlo.
\end{itemize}

Las pruebas a realizar de todo esto y mucho más se va a realizar en la lectura que sigue. Se recomienda leer detenidamente cada paso, y si no se entiende, releerlo. 
A eso sumarle el \textbf{pensamiento} de lograr la optimalidad, hasta acá se dijo sutilmente como es éste y su gran significado en las Ciencias de la Computación. A continuación, se explica exáctamente eso: en dónde se encasillan los \textbf{algoritmos polinomiales}, si es posible \textbf{verificar} una posible solución del problema en tiempo polinomial, y la relación de todo esto con los algoritmos no polinomiales. 

¿Es posible resolver el Problema del Juego de las Piedras, si yo \textbf{sé} resolver otro problema \textit{muy parecido} a éste?


\newpage
\section{Teoría de la Complejidad Computacional}

``A veces la respuesta es tan simple que da miedo, pero encontrarla es un infierno''. - Agente K (Hombres de Negro, 1997).

Planteamos formalmente el Problema del Juego de las Piedras en su versión de decisión: \textit{Dado un tablero de $n \times m$ celdas donde cada celda está vacía u ocupada por una piedra verde o celeste, ¿es posible sacar piedras de tal forma que cada fila tenga al menos una piedra y cada columna no tenga piedras de ambos colores?}

Si alguien nos da una posible solución $S$ (el tablero resultante despues de remover las piedras), a esa propuesta la llamamos \textit{certificado}. A priori no sabemos la correctitud de éste mismo, pero podemos verificarla con un algoritmo validador: el \textit{certificador}. Para asi verificar si cumple con todas las reglas.

Comunmente en los algoritmos implementados analizamos la complejidad temporal y espacial para medir cuánto tarda un algoritmo y cuánta memoria consume. Ahora, cuando pasamos de algoritmos polinómicos a algoritmos exponenciales, la barrera del tiempo se vuelve intratable (Tabla \ref{table: tabla1}).

\begin{table}[H]
    \centering
    \caption{Tiempos de ejecución aproximados de distintos órdenes de complejidad según el tamaño de entrada $n$, asumiendo un procesador que ejecuta un millón de instrucciones de alto nivel por segundo.}
    \renewcommand{\arraystretch}{1.3}
    \begin{tabular}{|l|l|l|l|l|l|l|l|}
    \hline
        Tamaño $n$ & $n$ & $n\log_{2}{n}$ & $n^{2}$ & $n^{3}$ & $1.5^{n}$ & $2^{n}$ & $n!$ \\ \hline
        $n = 10$ & $< 1$ seg & $< 1$ seg & $< 1$ seg & $< 1$ seg & $< 1$ seg & $< 1$ seg & 4 seg \\ \hline
        $n = 30$ & $< 1$ seg & $< 1$ seg & $< 1$ seg & $< 1$ seg & $< 1$ seg & 18 min & $10^{25}$ años \\ \hline
        $n = 50$ & $< 1$ seg & $< 1$ seg & $< 1$ seg & $< 1$ seg & 11 min & 36 años & muy largo \\ \hline
        $n = 100$ & $< 1$ seg & $< 1$ seg & $< 1$ seg & 1 seg & 12.892 años & $10^{17}$ años & muy largo \\ \hline
        $n = 1.000$ & $< 1$ seg & $< 1$ seg & 1 seg & 18 min & muy largo & muy largo & muy largo \\ \hline
        $n = 10.000$ & $< 1$ seg & $< 1$ seg & 2 min & 12 días & muy largo & muy largo & muy largo \\ \hline
        $n = 100.000$ & $< 1$ seg & 2 seg & 3 horas & 32 años & muy largo & muy largo & muy largo \\ \hline
        $n = 1.000.000$ & 1 seg & 20 seg & 12 días & 31.710 años & muy largo & muy largo & muy largo \\ \hline
    \end{tabular}
    \label{table:tabla1}
\end{table}

\subsection{Polinomialidad, NP, NP-Completitud y Reducciones}

Los problemas de decisión que pueden resolverse mediante un algoritmo en tiempo polinomial pertenecen al conjunto $\mathcal{P}$. Por otro lado, todo problema de decisión para el cual, dada una solución candidata (certificado), se puede verificar su validez en tiempo polinomial mediante un certificador eficiente, pertenece al conjunto $\mathcal{NP}$. Y nos podemos dar cuenta que, $\mathcal{P} \subseteq \mathcal{NP}$: si podemos resolver un problema desde cero en tiempo polinomial, también podemos verificar una solución ignorando el certificado y resolviéndolo nosotros mismos en tiempo polinomial.

En las Ciencias de la Computacion, hay una pregunta que hasta el día de hoy no se pudo responder: ¿se cumple que $\mathcal{P} = \mathcal{NP}$? Lo que muchos piensan (y lo de sentido común) es que $\mathcal{P} \neq \mathcal{NP}$. Que verificar la validez de una solución es mas fácil que encontrarla desde cero. Si llegara a demostrarse que $\mathcal{P} = \mathcal{NP}$, el impacto en el mundo seria tremendo, desde una optimización perfecta hasta la ruptura inmediata de toda la criptografia moderna.

Dentro de este universo definimos dos clases fundamentales:
\begin{itemize}
    \item Un problema $X$ es \textbf{NP-Hard} si todo problema $Y \in \mathcal{NP}$ puede reducirse polinomialmente a $X$ ($\forall Y \in \mathcal{NP}, Y \leq_P X$). Es decir, $X$ es \textit{al menos tan difícil} como cualquier problema de $\mathcal{NP}$.
    \item Un problema $X$ es \textbf{NP-Completo} si cumple al mismo tiempo dos condiciones: $X \in \mathcal{NP}$ y $X \in \text{NP-Hard}$.
\end{itemize}

Para demostrar que un problema nuevo $X$ es NP-Completo sin tener que reducir todos los problemas de $\mathcal{NP}$ uno por uno, usamos la transitividad de las \textbf{reducciones polinomiales} ($\leq_P$). Si agarramos un problema $Y$ que ya sabemos que es NP-Completo y demostramos que $Y \leq_P X$, probamos que si existiera una ``caja negra'' capaz de resolver $X$ en tiempo polinomial, podríamos transformar cualquier instancia de $Y$ en una instancia de $X$ en tiempo polinomial y resolver $Y$ eficientemente. Por lo tanto, $X$ es al menos tan difícil como $Y$, concluyendo que $X$ es NP-Completo.

\subsection{Pertenencia a NP}

Para demostrar que el Problema del Juego de las Piedras está en $\mathcal{NP}$, proponemos un algoritmo certificador eficiente (eficiente=polinomial) que recibe la instancia original del problema (\texttt{tablero\_original} de dimensiones $n \times m$) y el certificado propuesto (\texttt{tablero\_solucion} de dimensiones $n \times m$).

\begin{lstlisting}[mathescape=true, language=Python]
def validador_decision(tablero_original, tablero_solucion):
    n = len(tablero_original)
    if n == 0 or len(tablero_solucion) != n:
        return False
    m = len(tablero_original[0])
    
    # verificar subconjunto valido 
    for i in range(n):
        if len(tablero_solucion[i]) != m:
            return False
        for j in range(m):
            orig = tablero_original[i][j]
            sol = tablero_solucion[i][j]
            if sol not in ('.', 'V', 'C'):
                return False
            if sol != '.' and sol != orig:
                return False
    
    # verificar restriccion de columnas (a lo sumo tiene que haber un color por columna)
    for j in range(m):
        tiene_verde = False
        tiene_celeste = False
        for i in range(n):
            if tablero_solucion[i][j] == 'V':
                tiene_verde = True
            elif tablero_solucion[i][j] == 'C':
                tiene_celeste = True
            if tiene_verde and tiene_celeste:
                return False
    
                
    #verifico restriccion de filas (al menos una piedra por fila)
    for i in range(n):
        tiene_piedra = False
        for j in range(m):
            if tablero_solucion[i][j] in ('V', 'C'):
                tiene_piedra = True
                break
        if not tiene_piedra:
            return False
    
    return True
\end{lstlisting}

\textbf{Justificación de la Complejidad del Validador:}
\begin{enumerate}
    \item \textbf{Verificacion de subconjunto válido:} Se recorren las $n$ filas y por cada una, sus $m$ columnas realizando comparaciones de caracteres en tiempo constante $\mathcal{O}(1)$. El costo es $\mathcal{O}(n \cdot m)$.
    \item \textbf{Verificacion de colores por columna:} Se iteran las $m$ columnas y para cada una se inspeccionan sus $n$ celdas manteniendo dos banderas booleanas (\texttt{tiene\_verde} y \texttt{tiene\_celeste}). Cada celda se evalúa en $\mathcal{O}(1)$, resultando en $\mathcal{O}(n \cdot m)$.
    \item \textbf{Verificación de cobertura por fila:} Se recorren las $n$ filas verificando en a lo sumo $m$ pasos si existe al menos una piedra restante. El costo es $\mathcal{O}(n \cdot m)$.
\end{enumerate}

Sumando todo, la complejidad temporal total del certificador es $\mathcal{O}(n \cdot m) + \mathcal{O}(n \cdot m) + \mathcal{O}(n \cdot m) = \mathcal{O}(n \cdot m)$, que es lineal respecto al tamaño de la entrada (la cantidad de celdas del tablero) y estrictamente polinomial. La complejidad espacial adicional es $\mathcal{O}(1)$. Queda demostrado que \textbf{el Problema del Juego de las Piedras pertenece a $\mathcal{NP}$}.

\subsection{Demostración de NP-Completitud}

Habiendo probado que el problema del Juego de las Piedras $\in \mathcal{NP}$, vamos a demostrar que es NP-Completo reduciendo polinomialmente el problema \textbf{3-SAT} (conocido problema NP-Completo) a nuestro problema. Queremos:
\[ \text{3-SAT} \leq_P \text{Juego de las Piedras} \]

\noindent\textbf{Definición de 3-SAT:} Dados un conjunto de $m$ variables booleanas $X = \{x_1, x_2, \dots, x_m\}$ y un conjunto de $n$ clausulas $C = \{C_1, C_2, \dots, C_n\}$, donde cada cláusula $C_i = (\ell_{i,1} \lor \ell_{i,2} \lor \ell_{i,3})$ con 3 literales\footnote{Una literal es para decir que la variable puede ser $x_j$ o su negación $\neg x_j$}, ¿existe una asignación de valores de verdad para las variables tal que todas las cláusulas se evalúen como verdaderas?

\noindent\textbf{Construcción de la Transformación Polinomial:} \\
Dada una instancia arbitraria de 3-SAT con $n$ cláusulas y $m$ variables, construimos una instancia del Juego de las Piedras en un tablero $T$ de dimensiones $n \times m$ (donde cada fila $i \in \{1, \dots, n\}$ representa la cláusula $C_i$, y cada columna $j \in \{1, \dots, m\}$ representa la variable $x_j$) aplicando la siguiente regla para cada celda $T[i][j]$:
\begin{itemize}
    \item Si la variable $x_j$ aparece de forma \textbf{normal} ($x_j$) en la cláusula $C_i$, ponemos una piedra \textbf{Verde (\texttt{V})} en la celda $T[i][j]$.
    \item Si la variable $x_j$ aparece \textbf{negada} ($\neg x_j$) en la cláusula $C_i$, ponemos una piedra \textbf{Celeste (\texttt{C})} en la celda $T[i][j]$.
    \item Si la variable $x_j$ \textbf{no aparece} en la cláusula $C_i$, dejamos la celda $T[i][j]$ \textbf{vacía (\texttt{.})}.
\end{itemize}
(Nota: Si una clausula tuviera al mismo tiempo $x_j \lor \neg x_j$, esa cláusula es verdadera y se descarta en el preprocesamiento)

Por ejemplo, dada la fórmula $\mathcal{L} = (x_1 \lor \neg x_2 \lor x_3) \land (\neg x_1 \lor x_2 \lor \neg x_4) \land (\neg x_2 \lor \neg x_3 \lor x_4)$, la transformacion genera el tablero $3 \times 4$ de la Figura~\ref{fig:reduccion_3sat}.

\begin{figure}[H]
    \centering
    \begin{tikzpicture}[
    scale=0.85,
    transform shape,
    celda/.style={draw=gray!60, fill=gray!10, line width=0.6pt},
    verde/.style={circle, fill=green!60!black, draw=green!40!black, minimum size=0.62cm, inner sep=0pt, font=\footnotesize\sffamily\bfseries\color{white}},
    celeste/.style={circle, fill=cyan!75!blue!80, draw=cyan!50!black, minimum size=0.62cm, inner sep=0pt, font=\footnotesize\sffamily\bfseries\color{white}}
]
    \draw[celda] (0,0) rectangle (4,3);
    \draw[celda] (0,0) grid (4,3);

    % Encabezados de columnas (Variables)
    \node[font=\small\bfseries] at (0.5, 3.35) {$x_1$};
    \node[font=\small\bfseries] at (1.5, 3.35) {$x_2$};
    \node[font=\small\bfseries] at (2.5, 3.35) {$x_3$};
    \node[font=\small\bfseries] at (3.5, 3.35) {$x_4$};

    % Etiquetas de filas (Cláusulas)
    \node[anchor=east, font=\small] at (-0.2, 2.5) {$C_1 = (x_1 \lor \neg x_2 \lor x_3)$};
    \node[anchor=east, font=\small] at (-0.2, 1.5) {$C_2 = (\neg x_1 \lor x_2 \lor \neg x_4)$};
    \node[anchor=east, font=\small] at (-0.2, 0.5) {$C_3 = (\neg x_2 \lor \neg x_3 \lor x_4)$};

    % Piedras Fila 1
    \node[verde] at (0.5, 2.5) {V};
    \node[celeste] at (1.5, 2.5) {C};
    \node[verde] at (2.5, 2.5) {V};

    % Piedras Fila 2
    \node[celeste] at (0.5, 1.5) {C};
    \node[verde] at (1.5, 1.5) {V};
    \node[celeste] at (3.5, 1.5) {C};

    % Piedras Fila 3
    \node[celeste] at (1.5, 0.5) {C};
    \node[celeste] at (2.5, 0.5) {C};
    \node[verde] at (3.5, 0.5) {V};
\end{tikzpicture}


    \renewcommand{\thefigure}{3}
    \caption{Transformación polinomial de una instancia de 3-SAT a un tablero del Juego de las Piedras.}
    \label{fig:reduccion_3sat}
\end{figure}

Crear y llenar esta grilla de $n \times m$ leyendo las $n$ clausulas de $3$ literales cuesta $\mathcal{O}(n \cdot m)$ operaciones, por lo que \textbf{la reducción se ejecuta en tiempo polinomial}.

\noindent\textbf{Demostración de Equivalencia ($\iff$):} \\
Tenemos que probar que la formula 3-SAT $\mathcal{L}$ (la que tiene todo) es satisfacible \textbf{si y sólo si} el tablero $T$ construido tiene solución en el Juego de las Piedras.

\begin{itemize}
    \item \textbf{Ida ($\Rightarrow$): Si 3-SAT es satisfacible, entonces el tablero $T$ tiene solución.} \\
    Supongamos que existe una asignacion de verdad valida para $\mathcal{L}$. Para cada columna $j \in \{1, \dots, m\}$, hago en el tablero $T$:
    \begin{itemize}
        \item Si $x_j = \text{True}$, saco todas las piedras Celestes (\texttt{C}) de la columna $j$, quedandome solo con las Verdes (\texttt{V}).
        \item Si $x_j = \text{False}$, saco todas las piedras Verdes (\texttt{V}) de la columna $j$, quedandome solo con las Celestes (\texttt{C}).
    \end{itemize}
    Por la construccion, no va a haber columna que tenga piedras de los dos colores (se cumple la Regla 2 del problema). Además, como la asignación satisface $\mathcal{L}$, toda cláusula $C_i$ tiene al menos un literal verdadero. Si ese literal es positivo ($x_j = \text{True}$), en $T[i][j]$ había una piedra Verde que se quedó; si el literal es negado ($\neg x_j$ con $x_j = \text{False}$), en $T[i][j]$ había una piedra Celeste que también se quedó. Entonces, toda fila $i$ se queda con al menos una piedra (se cumple la Regla 1), obteniendo una solución válida para el Juego de las Piedras.

    \item \textbf{Vuelta ($\Leftarrow$): Si el tablero $T$ tiene solución, entonces 3-SAT es satisfacible.} \\
    Ahora suponemos que existe un conjunto de ``removidas'' valido en $T$ que cumple las dos reglas. Como ninguna columna $j$ tiene piedras de los dos colores en el tablero que queda, en cada columna $j$ quedan solo piedras Verdes, solo piedras celestes, o ninguna piedra. Definimos entonces la asignación booleana para cada variable $x_j$:
    \begin{itemize}
        \item Si en la columna $j$ quedaron piedras Verdes, ponemos $x_j = \text{True}$.
        \item Si en la columna $j$ quedaron piedras Celestes, ponemos $x_j = \text{False}$.
        \item Si la columna $j$ quedó completamente vacía, ponemos $x_j = \text{True}$.
    \end{itemize}
    Como la solución del tablero cumple la Regla 1, toda fila $i$ conserva al menos una piedra en alguna columna $j$. Si la piedra es Verde, por construcción el literal $x_j$ pertenecía a $C_i$ y nuestra asignación fijó $x_j = \text{True}$, haciendo verdadera a $C_i$. Si la piedra conservada es Celeste, el literal $\neg x_j$ pertenecía a $C_i$ y nuestra asignación fijó $x_j = \text{False}$, por lo que $\neg x_j = \text{True}$ y la cláusula $C_i$ también resulta verdadera. Así, las $n$ cláusulas se satisfacen simultáneamente.
\end{itemize}

Queda demostrado que $\text{3-SAT} \leq_P \text{Juego de las Piedras}$. Como el Juego de las Piedras $\in \mathcal{NP}$ y 3-SAT es NP-Completo, concluimos que \textbf{el Problema del Juego de las Piedras es NP-Completo}.








\newpage
\section{Búsqueda de la Exactitud: Backtracking y PLE}

je


\subsection{Teoría Backtracking y Podas}


\subsection{Diseño e Implementación del Backtracking}
        . lógica, podas,
        . codigo 
        . analisis complejidad (temporal y espacial)


\label{rec:codigo} % para llamarlo después. en el analisis de la complejidad.

\begin{lstlisting}[mathescape=true, language=Python]
def

\end{lstlisting}


\subsection{Modelo de Programación Lineal Entera}
. concepto PL + PLC/PLE y por que el PLE
. definicion variables + constantes/datos + funcion objetivo + restricciones y total de inecuaciones

\vfill
\section{Aproximaciones}

jeje

\subsection{Teoría}
. trade-off entre tiempo polinomial y optimalidad garantizada (elk r(A))


\subsection{Análisis del Algoritmo del Gringo Hinz}
. Demostración de validez, complejidad O(N.M)
. demostracion matemática de la cota r(A)


\subsubsection{Propuesta de Aproximación Propia}
. el por qué greedy + diseño determinístico
. pseudocodigo
. justificacion de la complejidad


\subsection{Pensamiento Algorítmico en la Vida Real}
. Conexión entre los 4 temas (Herbert Simon y la Racionalidad Limitada)
    . saber cuándo podar un camino con Backtracking y cuándo preferir una aproximación a tiempo frente a la parálisis por análisis).

    
tengo pensado:
En la vida cotidiana y en la ingeniería nos topamos con problemas "NP-Completos" (decisiones con miles de variables interconectadas donde buscar la perfección absoluta toma tiempo infinito). 
Ahí aplicamos los 4 pilares: Reducciones (reconocer que un problema nuevo es en realidad uno viejo disfrazado), Backtracking (explorar caminos pero tener la madurez de podar y retroceder apenas vemos que una rama no supera lo que ya logramos), Programación Lineal (optimizar nuestros objetivos dentro de restricciones rígidas), y Aproximaciones (el principio de Racionalidad Limitada o Satisficing del Premio Nobel y Premio Turing Herbert Simon: cuando encontrar el óptimo exacto es computacionalmente intratable, una buena aproximación a tiempo es infinitamente mejor que la perfección que llega tarde).

\newpage
\section{Experimentos y Mediciones}

. el por qué

\subsection{Generación de sets de datos}
. cómo se armaron los tableros de prueba (variando n, m, y densidad)


\subsection{Tiempos de ejecución}
. BT vs PL vs APROXIMACIONES


\subsection{Corroboración de la Cota r(A) y Calidad}
. Graficos comparando calidad del óptimo vs Gringo Hinz vs Greedy propio



\newpage
\section{Conclusión}

El desarrollo de este Trabajo Práctico nos permitio recorrer de punta a punta el ciclo de vida de un problema \textbf{NP-Completo}. A través de la reducción polinomial desde \textit{3-SAT}, demostramos formalmente por qué el problema que parecia sencillo, el juego de las piedras de Sophia esconde en sus reglas la misma dificultad que los problemas mas difíciles de la informática.

Al proponer su resolución exacta por \textbf{Backtracking} y \textbf{Programación Lineal}, comprobamos el valor de un buen modelado matemático: transformar la pregunta de ``qué piedras individuales borrar en una grilla de $n \times m$'' en ``qué color asignar a cada una de las $m$ columnas'' redujo drásticamente el espacio de búsqueda y permitió formular un modelo lineal limpio de $m + n$ restricciones. Igualmente, las mediciones confirmaron que el precio de exigir la perfección es el crecimiento exponencial del tiempo de la computadora.

Finalmente, el estudio de los \textbf{Algoritmos de Aproximación} demostró cómo es posible romper esa barrera exponencial ejecutando en tiempo lineal $\mathcal{O}(n \cdot m)$. Mientras que la estrategia aleatoria de ``El Gringo Hinz'' nos permitió demostrar analíticamente cómo incluso una moneda al aire garantiza una cota matemática sobre el valor esperado, el diseño de nuestra heurística \textbf{Greedy por Pricing Method} dejó una enseñanza fundamental: cuando combinamos la velocidad polinomial de una estrategia bien avariciosa con una medida inteligente de prioridad sobre las restricciones más frágiles, llegamos a lograr soluciones prácticamente optimas en milisegundos para volúmenes de datos que un algoritmo exacto tardaría siglos en resolver.



\vfill %para que la bibliografia vaya para abajo.
\printbibliography[heading=bibintoc, title={Bibliografía}]

\newpage
\end{document}